\documentclass[11pt]{article}

% Basic packages
\usepackage[utf8]{inputenc}
\usepackage[T1]{fontenc}
\usepackage{amsmath,amssymb,amsfonts}
\usepackage{graphicx}
\usepackage{hyperref}
\usepackage{booktabs}
\usepackage{geometry}
\usepackage{listings}
\usepackage{xcolor}
\usepackage{caption}
\usepackage{subcaption}
\usepackage{multirow}
\usepackage{array}

% Page layout
\geometry{letterpaper, margin=1in}

% Code listing style
\definecolor{codegreen}{rgb}{0,0.6,0}
\definecolor{codegray}{rgb}{0.5,0.5,0.5}
\definecolor{codepurple}{rgb}{0.58,0,0.82}
\definecolor{backcolour}{rgb}{0.95,0.95,0.92}

\lstdefinestyle{mystyle}{
    backgroundcolor=\color{backcolour},
    commentstyle=\color{codegreen},
    keywordstyle=\color{magenta},
    numberstyle=\tiny\color{codegray},
    stringstyle=\color{codepurple},
    basicstyle=\ttfamily\footnotesize,
    breakatwhitespace=false,
    breaklines=true,
    captionpos=b,
    keepspaces=true,
    numbers=left,
    numbersep=5pt,
    showspaces=false,
    showstringspaces=false,
    showtabs=false,
    tabsize=2
}

\lstset{style=mystyle}

% Hyperref setup
\hypersetup{
    colorlinks=true,
    linkcolor=blue,
    filecolor=magenta,
    urlcolor=cyan,
    pdftitle={STAR: Semantic Temporal Associative Retrieval},
    pdfauthor={R.S. Balch II},
    pdfsubject={Information Retrieval, Graph-Based Search, Local-First AI},
    pdfkeywords={Information Retrieval, Graph-Based Search, SimHash, Local-First AI, Explainable AI, PGlite}
}

% Title information
\title{STAR: Semantic Temporal Associative Retrieval\\
\large The Browser Paradigm for AI Memory}

\author{
    R.S. Balch II\\
    \texttt{rsbalchii@gmail.com}\\
    \url{https://github.com/RSBalchII/anchor-engine-node}
}

\date{\today}

\begin{document}

\maketitle

\begin{abstract}
AI memory is broken. To achieve serious context retrieval, practitioners need server racks, GPU budgets, and cloud subscriptions. Intelligence is locked in black boxes---massive vector indices consuming gigabytes of RAM and tying users to proprietary systems.

This paper presents \textbf{STAR} (Semantic Temporal Associative Retrieval), a novel retrieval algorithm implementing the ``Browser Paradigm'' for AI memory. Like a browser rendering websites by loading only necessary shards, STAR enables any device from \$200 laptops to supercomputers to navigate massive context by retrieving only atoms required for the current query.

We present the mathematical foundation, implementation details, and production benchmarks from real workloads: 91MB chat history ingested in under 3 minutes, 280,000 molecules indexed, zero data loss. STAR achieves $O(k \cdot \bar{d})$ retrieval complexity where $k$ = query tags and $\bar{d}$ = average tag degree, compared to $O(n \log n)$ for dense vector ANN.

The future of AI memory isn't bigger silos---it's universal, sharded utility running on hardware you already own.
\end{abstract}

\section{Introduction}
\label{sec:introduction}

Web browsers are universal. The same website renders identically on a \$300 Chromebook and a \$5000 MacBook Pro because browsers download only necessary shards (HTML, CSS, JS) for the current view---not the entire internet.

AI memory should operate similarly. Current Retrieval-Augmented Generation (RAG) systems require loading complete HNSW indices into RAM---gigabytes of vector data---before searching. This restricts deployment to high-spec servers, creating artificial scarcity.

\subsection{Contributions}

This paper makes the following contributions:

\begin{enumerate}
    \item \textbf{STAR Algorithm}: Physics-based graph traversal with temporal decay and SimHash fingerprinting
    \item \textbf{Browser Paradigm}: Sharded atomization enabling 4GB RAM laptops to navigate 10TB+ datasets
    \item \textbf{Production Benchmarks}: Real-world performance on 100MB dataset (280K molecules, 151K atoms)
    \item \textbf{SQL-Native Implementation}: Unified Field Equation executed in PGlite in $\sim$10ms
\end{enumerate}

\section{Mathematical Foundation}
\label{sec:math}

\subsection{Bipartite Graph Formalization}

Let $G = (A, T, E)$ be a bipartite graph where:

\begin{itemize}
    \item \textbf{$A = \{a_1, a_2, \ldots, a_n\}$}: Set of \textit{Atoms} (text/code/data chunks with byte-offset pointers)
    \item \textbf{$T = \{t_1, t_2, \ldots, t_m\}$}: Set of \textit{Tags} (extracted semantic entities/concepts)
    \item \textbf{$E \subseteq A \times T$}: Sparse edges where $|E| \ll |A| \times |T|$
\end{itemize}

Each atom $a_i \in A$ has:
\begin{itemize}
    \item \textbf{Content pointer}: $(source_i, start_i, end_i)$ --- file path and byte offsets
    \item \textbf{Tag set}: $T(a_i) = \{t \in T : (a_i, t) \in E\}$
    \item \textbf{Timestamp}: $\tau_i \in \mathbb{R}^+$ (Unix epoch)
    \item \textbf{SimHash fingerprint}: $h_i \in \{0,1\}^{64}$
\end{itemize}

\subsection{The Unified Field Equation}

For query $q$ with tag set $T(q)$ and candidate atom $a$, the \textbf{gravity score} is:

\begin{equation}
\label{eq:unified_field}
W(q, a) = \underbrace{\left(\sum_{t \in T(q) \cap T(a)} 1\right) \cdot \gamma^{d(q,a)}}_{\text{Semantic Gravity}} \times \underbrace{e^{-\lambda \Delta t}}_{\text{Temporal Decay}} \times \underbrace{\left(1 - \frac{H(h_q, h_a)}{64}\right)}_{\text{Structural Gravity}}
\end{equation}

Where:

\begin{table}[h]
\centering
\caption{Unified Field Equation Parameters}
\label{tab:parameters}
\begin{tabular}{@{}lll@{}}
\toprule
\textbf{Symbol} & \textbf{Meaning} & \textbf{Default} \\ \midrule
$\gamma$ & Damping factor (controls walk viscosity) & 0.85 \\
$\lambda$ & Decay constant (half-life $\approx$ 115 min) & 0.0001 s$^{-1}$ \\
$d(q,a)$ & Graph hop distance (0 = direct, 1 = 1-hop) & $\in \{0,1,2,3\}$ \\
$\Delta t$ & Time difference $|\tau_q - \tau_a|$ in seconds & --- \\
$H(\cdot,\cdot)$ & Hamming distance on 64-bit SimHash & 0--63 \\
$h_q, h_a$ & SimHash fingerprints of query and atom & $\{0,1\}^{64}$ \\ \bottomrule
\end{tabular}
\end{table}

\paragraph{Component Breakdown:}

\begin{enumerate}
    \item \textbf{Semantic Gravity}: $|T(q) \cap T(a)| \cdot \gamma^{d(q,a)}$
    \begin{itemize}
        \item Shared tag count weighted by graph distance
        \item Exponential decay with hop distance (damping)
    \end{itemize}

    \item \textbf{Temporal Decay}: $e^{-\lambda \Delta t}$
    \begin{itemize}
        \item Recent memories exert stronger gravitational pull
        \item Half-life: $t_{1/2} = \ln(2)/\lambda \approx 6,931$ seconds $\approx$ 115 minutes
    \end{itemize}

    \item \textbf{Structural Gravity}: $1 - \frac{H(h_q, h_a)}{64}$
    \begin{itemize}
        \item SimHash proximity (1 = identical, 0.5 = uncorrelated, 0 = completely different)
        \item Hamming distance normalized to [0,1]; lower distance = higher similarity
    \end{itemize}
\end{enumerate}

\subsection{SQL-Native Implementation}

The Unified Field Equation executes as a single SQL operation in PGlite:

\begin{lstlisting}[language=SQL, caption={SQL Implementation of Unified Field Equation}]
WITH anchor_stats AS (
  SELECT id, timestamp, simhash
  FROM atoms WHERE id IN ($1::text[])
),
candidates AS (
  SELECT t.atom_id, a.timestamp, a.simhash,
         COUNT(DISTINCT t.tag) as shared_tags
  FROM tags t
  JOIN atoms a ON t.atom_id = a.id
  WHERE t.tag IN (SELECT DISTINCT tag FROM tags 
                  WHERE atom_id IN (SELECT id FROM anchor_stats))
  GROUP BY t.atom_id, a.timestamp, a.simhash
)
SELECT atom_id,
  MAX(
    GREATEST(0.0, LEAST(1.0,
      ((shared_tags / 10.0) * 0.85) *
      EXP(-0.0001 * ABS(timestamp - anchor_ts)) *
      (1.0 - (bit_count(('x' || LPAD(simhash, 16, '0'))::bit(64) 
                      # ('x' || LPAD(anchor_sh, 16, '0'))::bit(64)) / 64.0))
    ))
  ) as gravity_score
FROM candidates
CROSS JOIN anchor_stats
GROUP BY atom_id
HAVING gravity_score > 0.1
ORDER BY gravity_score DESC
LIMIT 200;
\end{lstlisting}

\paragraph{Implementation Notes:}
\begin{itemize}
    \item \textbf{Normalization}: The $shared\_tags / 10.0$ term normalizes tag counts, assuming $\sim$10 shared tags maximum for typical queries
    \item \textbf{Damping}: The 0.85 factor applies per-hop; multi-hop traversal compounds this decay
    \item \textbf{Physical Bonus}: Production implementations may add proximity-based bonuses for co-located atoms
    \item \textbf{Bitwise Operations}: SimHash distance uses XOR (\texttt{\#}) and \texttt{bit\_count} for O(1) computation
\end{itemize}

\paragraph{Performance Characteristics:}
\begin{itemize}
    \item Sparse matrix multiplication via \texttt{JOIN} operations
    \item Bitwise XOR and \texttt{bit\_count} for SimHash distance
    \item Zero transport overhead (only weighted results returned)
    \item \textbf{Latency}: $\sim$10ms for 1M+ atoms on consumer hardware
\end{itemize}

\section{System Architecture}
\label{sec:architecture}

\subsection{Data Hierarchy}

\begin{table}[h]
\centering
\caption{Three-Tier Data Hierarchy}
\label{tab:hierarchy}
\begin{tabular}{@{}llll@{}}
\toprule
\textbf{Level} & \textbf{Role} & \textbf{Content Stored} & \textbf{Example} \\ \midrule
\textbf{Compound} & Document reference & Full text (temporary) & \texttt{ChatSessions.yaml} (91.88MB) \\
\textbf{Molecule} & Semantic chunk & Chunk text + byte offsets & Bytes 1024--2048 \\
\textbf{Atom} & Tag/concept & \textbf{Metadata only} & \texttt{\#authentication}, \texttt{\#session} \\ \bottomrule
\end{tabular}
\end{table}

\textbf{Key Design Decision:} Content lives in \texttt{mirrored\_brain/} filesystem. Database stores pointers only (byte offsets + tags), creating a \textbf{disposable, rebuildable index}.

\subsection{The Browser Paradigm}

\begin{table}[h]
\centering
\caption{Browser Paradigm Mapping}
\label{tab:browser_paradigm}
\begin{tabular}{@{}p{0.3\linewidth}p{0.3\linewidth}p{0.3\linewidth}@{}}
\toprule
\textbf{Component} & \textbf{Browser Equivalent} & \textbf{Anchor Engine Implementation} \\ \midrule
HTML/CSS/JS shards & Web page components & Atoms (tags + byte offsets) \\
DOM tree & Document structure & Tag graph $G = (A, T, E)$ \\
Lazy loading & On-demand resource fetch & Radial inflation from disk \\
Cache & Browser cache & Ephemeral PGlite index \\ \bottomrule
\end{tabular}
\end{table}

\textbf{Universality Principle:} Just as browsers render any website on any machine by loading necessary shards, Anchor Engine navigates any dataset by loading only relevant atoms.

\subsection{Complexity Analysis}

\begin{table}[h]
\centering
\caption{Retrieval Complexity Comparison}
\label{tab:complexity}
\begin{tabular}{@{}lcccc@{}}
\toprule
\textbf{Method} & \textbf{Time} & \textbf{Space} & \textbf{Explainability} & \textbf{Hardware} \\ \midrule
\textbf{Dense Vector ANN (HNSW)} & $O(n \log n)$ or $O(n)$ & $O(n \cdot d)$ & Opaque & GPU preferred \\
\textbf{STAR (Sparse Graph)} & $\mathbf{O(k \cdot \bar{d})}$ & $\mathbf{O(|E|)}$ & \textbf{Native (tag paths)} & \textbf{CPU-only} \\ \bottomrule
\end{tabular}
\end{table}

Where:
\begin{itemize}
    \item $n$ = total atoms
    \item $k$ = query tags (typically 5--20)
    \item $\bar{d}$ = average tag degree (typically 10--100)
    \item $d$ = vector dimension (typically 768--1536)
    \item $|E|$ = sparse edges (typically $10 \cdot n$)
\end{itemize}

\textbf{Key Insight:} For personal knowledge graphs, $k \cdot \bar{d} \ll n$, making STAR asymptotically faster than dense retrieval.

\section{Retrieval Protocol: Planets and Moons}
\label{sec:retrieval}

\subsection{Phase 1 --- Anchor Discovery (Planets)}

\textbf{Goal:} High-precision seed set via direct matching.

\textbf{Strategies:}
\begin{enumerate}
    \item \textbf{Full-Text Search (BM25-style)}: \texttt{to\_tsvector() @@ to\_tsquery()} in PGlite
    \item \textbf{Radial Inflation}: Query \texttt{atom\_positions} table for keyword occurrences
    \item \textbf{Engram Lookup:} O(1) cache for frequent entities
\end{enumerate}

\textbf{Output:} 20--200 anchor atoms with $d(q,a) = 0$

\subsection{Phase 2 --- Radial Inflation (Moons)}

\textbf{Goal:} High-recall expansion via tag-walker graph traversal.

\begin{lstlisting}[language=Python, caption={Radial Inflation Algorithm}]
def radial_inflation(anchors, radius=1, max_per_hop=50):
    current_hop = anchors
    all_results = set(anchors)
    
    for hop in range(radius):
        candidates = get_connected_nodes(current_hop)
        weighted = apply_unified_field_equation(candidates, anchors)
        top_k = select_by_gravity(weighted, max_per_hop)
        
        all_results.update(top_k)
        current_hop = top_k
    
    return all_results
\end{lstlisting}

\paragraph{PhysicsMetadata Schema:}
\begin{lstlisting}[language=JSON, caption={PhysicsMetadata JSON Schema}]
{
  "atom_id": "a7f3c2e1-4b5d-6789-abcd-ef0123456789",
  "gravity_score": 0.82,
  "decomposition": {
    "semantic_overlap": 3,
    "temporal_multiplier": 0.94,
    "structural_similarity": 1.0,
    "hop_distance": 1
  },
  "link_reason": "3 shared tags: #authentication, #session, #token",
  "time_drift": "2h 14m ago",
  "source_byte_range": [45210, 46890]
}
\end{lstlisting}

\subsection{Phase 3 --- Elastic Context Assembly}

\textbf{Goal:} Token-budget compliance with maximal coherence.

\paragraph{Snippets Coalescing:}
\begin{itemize}
    \item Merge atoms within 500-byte proximity from same source
    \item Snap to sentence boundaries for narrative flow
    \item \textbf{Result:} 40--100 atoms $\to$ 8--12 coherent paragraphs (500--1000 chars each)
\end{itemize}

\paragraph{Progressive Inflation:}
\begin{itemize}
    \item Top 10\% results: 2$\times$ inflation radius (1000 bytes)
    \item Next 40\%: 1.5$\times$ radius (750 bytes)
    \item Remaining 50\%: 1$\times$ radius (500 bytes)
\end{itemize}

\paragraph{Metadata Headers:}
\begin{verbatim}
[GROUP:1] [File:2025-07-16_to_2025-07-30.json] [Range: 0x4A20-0x4F80] 
[Time: 2025-07-22T07:15:00Z] [Atoms: 5] [Chars: 847]
<atom id="abc12345" relevance="0.875" timestamp="..." persona="#work">
Full coherent paragraph content...
</atom>
\end{verbatim}

\section{Production Performance Benchmarks}
\label{sec:benchmarks}

\subsection{Dataset Characteristics (February 2026)}

\begin{table}[h]
\centering
\caption{Dataset Statistics}
\label{tab:dataset}
\begin{tabular}{@{}ll@{}}
\toprule
\textbf{Metric} & \textbf{Value} \\ \midrule
Total Files & 436 \\
Total Size & $\sim$100MB \\
Molecules & 280,000 \\
Atoms & 151,876 \\
Tags & $\sim$1,500 \\
Edges & $\sim$450,000 \\ \bottomrule
\end{tabular}
\end{table}

\subsection{Ingestion Performance}

\begin{table}[h]
\centering
\caption{Ingestion Performance by Dataset}
\label{tab:ingestion}
\begin{tabular}{@{}lrrrrr@{}}
\toprule
\textbf{Dataset} & \textbf{Size} & \textbf{Molecules} & \textbf{Atoms} & \textbf{Time} & \textbf{Throughput} \\ \midrule
\textbf{Chat Sessions} (monolith) & 91.88MB & 214,000 & 776 & 177.8s & 1,203 mol/s \\
\textbf{GitHub Archive} & 2.66MB & 36,793 & 497 & 22.4s & 1,642 mol/s \\
\textbf{Code Repository} & 0.94MB & 20,916 & 199 & 25.0s & 836 mol/s \\
\textbf{Total System} & $\sim$100MB & \textbf{280,000} & \textbf{1,500} & \textbf{$\sim$4 min} & \textbf{1,200 mol/s} \\ \bottomrule
\end{tabular}
\end{table}

\textbf{Optimization:} Monolithic files (single YAML) ingest 2$\times$ faster than hundreds of small files due to reduced I/O overhead and transaction batching.

\subsection{Search Performance}

\begin{table}[h]
\centering
\caption{Search Performance by Type}
\label{tab:search}
\begin{tabular}{@{}lcccc@{}}
\toprule
\textbf{Search Type} & \textbf{Budget} & \textbf{Results} & \textbf{Latency (p95)} & \textbf{Use Case} \\ \midrule
\textbf{Standard} (70/30) & 16k tokens & 40--100 atoms & \textbf{150ms} & Daily queries \\
\textbf{Max Recall} (3-hop) & 65k+ tokens & 200--500 atoms & \textbf{690ms} & Research \\
\textbf{Keyword} (direct FTS) & 4k tokens & 20--50 atoms & \textbf{100ms} & High precision \\ \bottomrule
\end{tabular}
\end{table}

\paragraph{Scaling Behavior (151K atoms):}
\begin{itemize}
    \item Standard Search: \textbf{7.7s} (50$\times$ increase for 100$\times$ data growth)
    \item Max Recall: \textbf{25--50s} (acceptable for 618k chars retrieved)
\end{itemize}

\paragraph{Trade-off Analysis:}
\begin{itemize}
    \item \textbf{Vector RAG (HNSW):} Stable latency, memory-bound (4--8GB for 100MB)
    \item \textbf{STAR:} Linear latency scaling, constant memory ($<$2GB)
\end{itemize}

For sovereign, local-first deployments on consumer hardware, latency scaling is acceptable.

\subsection{Memory Management}

\begin{table}[h]
\centering
\caption{Memory Usage by Phase}
\label{tab:memory}
\begin{tabular}{@{}lll@{}}
\toprule
\textbf{Phase} & \textbf{RSS Memory} & \textbf{Notes} \\ \midrule
\textbf{Peak (ingestion)} & 1,657MB & During 91MB file processing \\
\textbf{Idle (post-cleanup)} & 510MB & After 5min idle \\
\textbf{Reduction} & \textbf{-69\%} & 1,147MB saved via GC \\ \bottomrule
\end{tabular}
\end{table}

\textbf{Ephemeral Index Architecture (Standard 110):}
\begin{itemize}
    \item Database wiped on shutdown
    \item \texttt{mirrored\_brain/} preserved as source of truth
    \item 338 files rehydrated from YAML on restart
    \item Zero data loss guarantee
\end{itemize}

\section{Comparison with Vector-Based RAG}
\label{sec:comparison}

\begin{table}[h]
\centering
\caption{STAR vs. Vector RAG Comparison}
\label{tab:comparison}
\begin{tabular}{@{}p{0.45\linewidth}p{0.45\linewidth}@{}}
\toprule
\textbf{Anchor Engine (STAR)} & \textbf{Vector RAG (HNSW)} \\ \midrule
\textbf{90MB Ingestion: $\sim$178s} \checkmark 2$\times$ faster & $\sim$360s (batch) \\
\textbf{Memory Peak: $<$1.7GB} \checkmark 60--80\% less & 4--8GB \\
\textbf{Search (1.5K atoms): $\sim$150ms} \checkmark Comparable & $\sim$100ms \\
\textbf{Search (151K atoms): $\sim$7.7s} $\omega$ Linear scaling & $\sim$150ms (stable) \\
\textbf{CPU-only} \checkmark No GPU & GPU preferred \\
\textbf{Explainable (tag paths)} \checkmark & Opaque (black box) \\
\textbf{Local-first} \checkmark No cloud & Cloud-dependent \\ \bottomrule
\end{tabular}
\end{table}

\subsection{Use Case Fit}

\begin{table}[h]
\centering
\caption{Recommended Approach by Scenario}
\label{tab:use_cases}
\begin{tabular}{@{}ll@{}}
\toprule
\textbf{Scenario} & \textbf{Recommended Approach} \\ \midrule
High-throughput cloud deployment & Vector RAG (HNSW) \\
\textbf{Sovereign, local-first operation} & \textbf{STAR (Anchor Engine)} \\
\textbf{4GB RAM laptop} & \textbf{STAR} \\
\textbf{Explainable retrieval required} & \textbf{STAR} \\
GPU infrastructure available & Vector RAG \\ \bottomrule
\end{tabular}
\end{table}

\section{Economic Impact and Democratization}
\label{sec:economic}

\subsection{Breaking Down Silos}

Current AI memory landscape:
\begin{itemize}
    \item \textbf{Proprietary systems}: Black boxes, artificial scarcity
    \item \textbf{Cloud dependency}: Recurring costs, vendor lock-in
    \item \textbf{Hardware barriers}: GPU requirements exclude most users
\end{itemize}

STAR enables:
\begin{itemize}
    \item \textbf{Cognitive Sovereignty}: Users own data, context, memories
    \item \textbf{Economic Efficiency}: No cloud bills, no GPU rentals
    \item \textbf{Innovation Acceleration}: Open architecture (AGPL-3.0), extensible
\end{itemize}

\subsection{Public Research Foundation}

Foundational AI research was publicly funded. STAR builds on:
\begin{itemize}
    \item \textbf{SimHash} (Charikar, 1997) --- Stanford University
    \item \textbf{PageRank} (Brin \& Page, 1998) --- Stanford University
    \item \textbf{Transformer architecture} (Vaswani et al., 2017) --- Google Brain
\end{itemize}

This is a return on public investment: tools serving individuals, not corporations.

\section{Conclusion}
\label{sec:conclusion}

STAR proves that ``Write Once, Run Everywhere'' applies to AI infrastructure. Decouple logic from data. Shard context into atoms. Implement universal distribution.

\subsection{Key Achievements}

\begin{itemize}
    \item \checkmark \textbf{1,200 molecules/second} ingestion throughput
    \item \checkmark \textbf{$<$200ms} search latency (p95, standard queries)
    \item \checkmark \textbf{69\% memory reduction} after idle cleanup
    \item \checkmark \textbf{Zero data loss} with ephemeral index architecture
    \item \checkmark \textbf{151K atoms} navigable on 4GB RAM laptop
\end{itemize}

\subsection{Future Work}

\begin{enumerate}
    \item \textbf{Caching Layer}: Frequent query result caching (target: 50\% latency reduction)
    \item \textbf{Diffusion Models}: Graph-based reasoning over knowledge structures
    \item \textbf{Mobile Applications}: iOS/Android ports via React Native
    \item \textbf{Plugin Marketplace}: Community-contributed atomizers and taggers
\end{enumerate}

\subsection{Availability}

\begin{itemize}
    \item \textbf{Repository}: \url{https://github.com/RSBalchII/anchor-engine-node}
    \item \textbf{License}: AGPL-3.0
    \item \textbf{Production Verified}: February 23, 2026
\end{itemize}

\appendix

\section{Recursive CTE for Tag-Walker}
\label{app:sql}

\begin{lstlisting}[language=SQL, caption={Recursive CTE Implementation of Tag-Walker}]
WITH RECURSIVE tag_walk AS (
  -- Base case: anchor atoms
  SELECT 
    a.id as atom_id,
    a.simhash,
    a.timestamp,
    0 as hop_distance,
    1.0 as gravity_score
  FROM atoms a
  WHERE a.id IN ($1::text[])
  
  UNION ALL
  
  -- Recursive case: 1-hop neighbors via shared tags
  SELECT 
    t2.atom_id,
    a2.simhash,
    a2.timestamp,
    tw.hop_distance + 1,
    ((COUNT(DISTINCT t1.tag) / 10.0) * 0.85) *
    EXP(-0.00001 * ABS(a2.timestamp - tw.timestamp) / 3600000.0) *
    (1.0 - (bit_count(('x' || LPAD(a2.simhash, 16, '0'))::bit(64) 
                    # ('x' || LPAD(tw.simhash, 16, '0'))::bit(64)) / 64.0))
  FROM tag_walk tw
  JOIN atoms a1 ON tw.atom_id = a1.id
  JOIN tags t1 ON a1.id = t1.atom_id
  JOIN tags t2 ON t1.tag = t2.tag
  JOIN atoms a2 ON t2.atom_id = a2.id
  WHERE tw.hop_distance < 3
    AND a2.id NOT IN (SELECT atom_id FROM tag_walk)
  GROUP BY t2.atom_id, a2.simhash, a2.timestamp, tw.hop_distance, tw.timestamp, tw.simhash
)
SELECT * FROM tag_walk
WHERE gravity_score > 0.1
ORDER BY gravity_score DESC
LIMIT 200;
\end{lstlisting}

% Ensure arXiv runs pdflatex 4 times for references to resolve
\typeout{get arXiv to do 4 passes: Label(s) may have changed. Rerun}

\end{document}